\subsection{Restricted Dual and Matrix Coefficients}

PROPOSITION 3. --- Let \((\mathrm{V},\pi)\) be a finite-dimensional linear representation of the algebra A. The kernel \(\mathfrak{a}\) of \(\pi\) is a two-sided ideal of A of finite codimension, and \(\Theta_{\pi}(\mathrm{A})\) is the orthogonal of \(\mathfrak{a}\) in \(A^{*}\) . The transpose of the mapping \(\pi\) defines an isomorphism from the dual of the K-vector space \(\pi(\mathrm{A})\) to \(\Theta_{\pi}(\mathrm{A})\) .

Since \(\operatorname{End}_{\mathrm{K}}(\mathrm{V})\) is finite-dimensional over K, \(\mathfrak{a}\) is a two-sided ideal of A of finite codimension. The vector space \(\Theta_{\pi}(\mathrm{A})\) is a finite-dimensional subspace of \(A^{*}\) , and its orthogonal in A is equal to \(\mathfrak{a}\); by Theorem 7 of II, §7, No.~5, p.~301, the space \(\Theta_{\pi}(\mathrm{A})\) is therefore the orthogonal of \(\mathfrak{a}\) in \(A^{*}\) . Moreover, since \(\mathfrak{a}\) is the kernel of \(\pi\) , the transpose of the mapping \(\pi\) defines an isomorphism from the dual of \(\pi(\mathrm{A})\) to the orthogonal of \(\mathfrak{a}\) in \(A^{*}\) (II, §7, No.~5, p.~299, Corollary of Theorem 5).

COROLLARY. --- The restricted dual \(\Theta(\mathrm{A})\) of A is the set of coefficients of the finite-dimensional linear representations of A.

By definition, \(\Theta(A)\) is the union of the orthogonals of the two-sided ideals of \(A\) of finite codimension. We therefore have \(\Theta_{\pi}(A) \subset \Theta(A)\) for every finite-dimensional linear representation \(\pi\) of \(A\) . Conversely, let \(f\) be an element of \(\Theta(A)\) , and let \(\mathfrak{a}\) be a two-sided ideal of finite codimension, contained in the kernel of \(f\) (VIII, p.~376, Definition 2). Denote by \(\pi\) the linear representation of \(A\) in \(A/\mathfrak{a}\) deduced from the left regular representation in \(A\) by passing to the quotient. Let \(x\) be the class of 1 (mod \(\mathfrak{a}\) ), and let \(x^{*}\) be the linear form on \(A/\mathfrak{a}\) deduced from \(f\) . We have \(f = c_{\pi}(x, x^{*})\) , so that \(f\) is a coefficient of \(\pi\) .

\pandocbounded{\includegraphics[keepaspectratio]{images/part003_35b408b4.jpg}}

Take note that if \((V,\pi)\) is a linear representation of A that is not finite-dimensional over K, then the space \(\Theta_{\pi}(A)\) is not necessarily contained in \(\Theta(A)\) .

